% 問題の定式化

アニメ制作業界は、構造的な\textbf{貢献の不可視性問題（contribution invisibility problem）}を抱えている。
アニメを生み出す創造的労働の相当部分が記録されず、
帰属も示されず、検証もできない状態にある。
本節では、この問題を3つの側面から形式化し、
既存の解決策がなぜ不十分であるかを分析する。

% ----- 3.1 クレジットギャップ -----
\subsection{クレジットギャップ}
\label{subsec:credit_gap}

典型的なアニメ制作において、エンドクレジット（end credits）には
各話に貢献したスタッフの名前が記載される。
しかし、このクレジットシステムには3つの構造的限界がある。

\begin{enumerate}
  \item \textbf{下請けの不透明性：}
    作業が複数層の下請け構造を通じて外注される場合
    （業界における一般的な慣行。\cref{subsec:anime_workflow}参照）、
    二次下請けや三次下請けに所属する個々のアニメーターは
    クレジットから完全に省略されることが多い。
    下請けスタジオの名前のみが記載されるか、
    場合によっては一切のクレジット表記がなされないこともある。
  \item \textbf{枠の制約：}
    クレジットには実務上のスペースと時間の制約がある。
    テレビ放送ではクレジットに割り当てられる画面時間が限られており、
    制作委員会が編集上または契約上の理由から
    記載される貢献者の数を制限する場合もある。
  \item \textbf{スタッフの入れ替わり：}
    制作途中でプロジェクトを離脱したアニメーター---自発的な離脱であれ
    スケジュールの都合による離脱であれ---は、
    相当量の作業を完了していたとしても
    最終的なクレジットから名前が削除されることがある。
\end{enumerate}

これらの制約は、\textbf{クレジットが貢献の包括的な記録ではなく、
不完全かつ取捨選択された記録として機能している}ことを意味する。

% ----- 3.2 フリーランスアニメーターへの影響 -----
\subsection{フリーランスアニメーターへの影響}
\label{subsec:consequences}

クレジットギャップは、アニメーターの47--70\%を占める
フリーランサー~\cite{janica2023survey}に対して具体的な影響を及ぼす。

\begin{itemize}
  \item \textbf{ポートフォリオの非対称性：}
    新たな契約を獲得する際、フリーランサーは自身の実績を示す必要がある。
    クレジットがなければ、口頭での主張、個人的な人脈、
    あるいは独立した検証手段を欠く物理的なサンプルに頼らざるを得ない。
    フリーランサーを雇用するスタジオ側にも、
    候補者が主張する経歴を検証する標準的な方法がない。
  \item \textbf{交渉上の不利：}
    アニメーターは、クレジットされていない仕事を報酬交渉に活用できない。
    著名な作品で30カットを完成させた原画マンであっても、
    クレジットされていなければ、より高い報酬を正当化する
    検証可能な証拠を持たないことになる。
  \item \textbf{キャリアの脆弱性：}
    フリーランサーのキャリア履歴は、主に同僚の記憶や
    非公式な業界ネットワークの中にのみ存在する。
    スタジオが倒産した場合---アニメ業界では頻繁に起こることである---
    あるいは重要な人脈が引退した場合、
    これらの非公式な記録は消失する。
  \item \textbf{業界レベルでのデータ喪失：}
    集約的な観点からも、包括的な貢献記録の欠如は、
    労働力の動態を研究すること、特定の職種における人手不足を特定すること、
    あるいは効果的な産業政策を策定することを困難にしている。
    JAniCAやNAFCAによる調査は貴重な断面データ（スナップショット）を提供するが、
    プロジェクト単位の帰属データを捕捉することはできない。
  \item \textbf{社会参加への障壁：}
    その影響は業界内にとどまらない。
    住宅ローンの申請、ビザの手続き、銀行の信用評価には、
    いずれも安定した雇用と収入履歴の文書化された証明が求められる。
    検証可能なキャリア記録を持たないフリーランスアニメーターは、
    これらの手続きにおいて構造的な不利益を被る---技能や収入が
    不足しているからではなく、制度が要求する書類を
    提示できないためである。
    こうして貢献の不可視性問題は、アニメーターの
    生活上の選択（持ち家の取得、国際的な移動など）を
    業界自身からは見えない形で制約している。
\end{itemize}

% ----- 3.3 既存の解決策の限界 -----
\subsection{既存の解決策の限界}
\label{subsec:existing_limitations}

フリーランスアニメーターの問題に対処するいくつかの取り組みが近年行われているが、
貢献の不可視性問題を完全に解決するものはない。

\begin{itemize}
  \item \textbf{フリーランス保護法（Freelance Protection Act）（2024年）：}
    日本の新しい法律は、書面による契約の義務化や支払い期限の規制など、
    個人事業主に対する保護を導入している。
    しかし、貢献の記録や帰属の表示を義務付けてはおらず、
    法的に適正な契約を結んでいてもクレジットされないアニメーターが存在しうる。
  \item \textbf{NAFCA標準契約：}
    NAFCAは、スタジオとアニメーター間の契約テンプレートを策定している。
    これらは取引の透明性を向上させるが、
    二者間の文書であり、
    アニメーターが将来の仕事先に持ち運べる
    公開的に検証可能な記録を生成するものではない。
  \item \textbf{コミュニティデータベース（例：Sakugabooru、AniList）：}
    ファンが維持するデータベースがアニメーターの貢献を追跡しようと試みている。
    これらは有用であるが、ボランティアの努力に依存しており、
    本質的に不完全であり、公式な検証を欠いている。
  \item \textbf{集中型ポートフォリオプラットフォーム：}
    ArtStationや個人ウェブサイトなどのプラットフォームにより、
    アニメーターは自身の作品を展示できるが、
    暗号学的な証明のない自己申告に依存しており、
    利用規約を変更したりサービスを終了したりしうる
    単一のプラットフォーム運営者によって管理されている。
\end{itemize}

% ----- 3.4 問題の形式化 -----
\subsection{問題の形式化}
\label{subsec:formalization}

貢献の不可視性問題を以下のように定義する。
$A$ を制作物 $P$ に貢献したすべてのアニメーターの集合とする。
$C(P) \subseteq A$ を $P$ の公式記録においてクレジットされたアニメーターの集合とする。
\textbf{クレジットギャップ（credit gap）}は次のように表される。

\begin{equation}
  \Delta(P) = A \setminus C(P)
  \label{eq:credit_gap}
\end{equation}

任意のアニメーター $a \in \Delta(P)$ に対して、$P$ への貢献についての
検証可能かつ永続的かつ可搬性のある記録は存在しない。

理想的な解決策は以下の要件を満たすべきである。

\begin{enumerate}
  \item[\textbf{R1}] \textbf{完全性（Completeness）：} 下請け構造における位置に関係なく、
    すべてのアニメーター $a \in A$ が記録を取得できること。
  \item[\textbf{R2}] \textbf{検証可能性（Verifiability）：} 第三者が、元のスタジオに
    問い合わせることなく、記録の真正性を暗号学的に検証できること。
  \item[\textbf{R3}] \textbf{永続性（Persistence）：} スタジオの倒産、
    プラットフォームの停止、業界インフラの変更に対して記録が存続すること。
  \item[\textbf{R4}] \textbf{可搬性（Portability）：} アニメーターが記録を所有・管理し、
    将来の雇用主やクライアントに対して提示できること。
  \item[\textbf{R5}] \textbf{最小限の摩擦（Minimal friction）：} 記録プロセスが、アニメーターに
    追加の労力を要求することなく、既存のワークフローに統合されること。
\end{enumerate}

\Cref{tab:requirements}は、これらの要件に照らして既存の解決策を評価したものである。
AnimatorSBTは5つの要件すべてを満たすよう設計されている。

\begin{table}[ht]
\centering
\caption{5つの要件に対する既存の解決策の評価}
\label{tab:requirements}
\footnotesize
\setlength{\tabcolsep}{3pt}
\begin{tabular}{@{}lccccc@{}}
\toprule
\textbf{解決策} & \textbf{R1} & \textbf{R2} & \textbf{R3} & \textbf{R4} & \textbf{R5} \\
\midrule
公式クレジット         & --         & \checkmark & \checkmark & --         & \checkmark \\
フリーランス保護法     & --         & --         & --         & --         & \checkmark \\
NAFCA標準契約          & $\triangle$    & --         & --         & --         & \checkmark \\
コミュニティDB         & $\triangle$    & --         & $\triangle$    & --         & --         \\
ポートフォリオプラットフォーム & \checkmark & --         & --         & $\triangle$    & --         \\
\midrule
\textbf{AnimatorSBT}   & \checkmark & \checkmark & \checkmark & \checkmark & \checkmark \\
\bottomrule
\end{tabular}
\end{table}
